\begin{figure}[tbp]
\centering
\begin{tikzpicture}[
  >=Latex,
  line/.style={very thick},
  map/.style={->,thick},
  special/.style={circle,fill=black,inner sep=1.8pt},
  excluded/.style={circle,draw,fill=white,inner sep=2.1pt},
  every node/.style={font=\small}
]
  \node[anchor=east] at (-4.45,3.2) {$t$};
  \draw[line] (-4.2,3.2) -- (4.25,3.2);
  \foreach \x in {-3.55,-2.9,-2.25,-1.6,-0.95}
    \node[excluded] at (\x,3.2) {};
  \node[anchor=south west,font=\scriptsize,align=left] at (-3.65,3.35)
    {open circles: the five factors excluded from the regular locus};

  \node[anchor=east] at (-4.45,1.15) {$y$};
  \draw[line] (-4.2,1.15) -- (4.25,1.15);
  \draw[map] (-1.1,2.98) -- (-1.1,1.37)
    node[midway,left,align=right,font=\scriptsize]
      {$y=(t-1)^2/t$\\$2{:}1$};

  \node[anchor=east] at (-4.45,-1.15) {$z$};
  \draw[line] (-4.2,-1.15) -- (4.25,-1.15);
  \draw[map] (-1.1,0.93) -- (-1.1,-0.93)
    node[midway,right,align=left,font=\scriptsize]
      {$z=(y-y^{-1})^2/16$\\$4{:}1$};

  \node[excluded] at (-3.55,-1.15) {};
  \node[anchor=north,align=center,font=\scriptsize] at (-3.55,-1.3)
    {$-1/4$\\GRS boundary};
  \node[special] at (-0.55,-1.15) {};
  \node[anchor=north,align=center,font=\scriptsize] at (-0.55,-1.3)
    {$1$\\bracket collision};

  \node[draw,rounded corners,align=left,font=\scriptsize,inner sep=4pt,
    anchor=south west] at (0.1,1.45)
    {deck maps on $y$ and realizing permutations:\\
      $y\mapsto y$\quad id\qquad
      $y\mapsto-y$\quad $(5\,6)$\\
      $y\mapsto y^{-1}$\quad $(3\,5)(4\,6)$\qquad
      $y\mapsto-y^{-1}$\quad $(3\,5\,4\,6)$};
\end{tikzpicture}
\caption{The degree-eight quotient of the regular non-GRS pencil (special points on
the \(z\)-line are arranged schematically).  The parameter \(t\)
double-covers \(y=(t-1)^2/t\), and the four identifications
\(y\mapsto\pm y^{\pm1}\), each realized by a displayed coordinate permutation,
cover \(z\).  The open point \(z=-1/4\) is the excluded GRS boundary;
\(z=1\) is the bracket collision handled in the proof.}
\label{fig:pencil-quotient}
\end{figure}
